\chapter{Введение: зачем памяти нужен субстрат}

\section{Постановка задачи}

У локального AI-компаньона есть требование, которое почти невозможно
обойти архитектурной хитростью: он должен вести себя не как одноразовый
ответчик на очередной запрос, а как сущность с накопленной биографией. Это
означает по меньшей мере четыре вещи.

\begin{enumerate}
\item Он должен связывать события между собой, а не хранить их как изолированные записи.
\item Он должен ошибаться не хаотически, а осмысленно: делать предварительные выводы, сомневаться в них и позднее их пересматривать.
\item Он должен меняться во времени: недавнее должно звучать иначе, чем давнее; важное должно переживать неважное; повторяющееся должно сворачиваться в более устойчивые структуры.
\item Он должен вспоминать контекстно: не одинаково по любой команде, а в зависимости от текущего вопроса, эмоционального состояния и уже активированного окружения.
\end{enumerate}

Плоская история сообщений решает лишь малую часть этой задачи. Векторная база
воспоминаний продвигается дальше: она позволяет искать смыслово похожие записи.
Но ни история сообщений, ни набор независимых векторных элементов сами по себе
не дают ответа на более глубокий вопрос: где в памяти живут связи между фактами,
уверенность в них, их происхождение и их эволюция во времени?

\section{Почему здесь появляется граф}

Слово ``граф'' легко превратить в пустую модную вывеску. В этой книге под
графом понимается не престижная структура данных сама по себе, а минимально
естественный ответ на перечисленные выше требования.

Если память должна хранить не только ``что произошло'', но и ``с чем это
связано'', ``насколько этому доверяют'', ``из чего это было выведено'', ``как это
влияет на соседние убеждения'' и ``когда это было важно'', то у архитектуры
почти неизбежно возникают:

\begin{itemize}
\item узлы как носители содержания;
\item рёбра как носители отношений;
\item метаданные на узлах и рёбрах;
\item процессы над графом, а не просто чтение из хранилища.
\end{itemize}

\begin{intuition}[Граф здесь не украшение]
Если убрать из памяти отношения и оставить только список фрагментов текста,
система ещё сможет что-то ``вытаскивать по похожести'', но уже не сможет внятно
ответить, почему одно суждение усиливает другое, откуда возникла гипотеза, что
следует пересмотреть после опровержения факта и по какому маршруту вообще
течёт активация. Иначе говоря, в такой системе не останется самой механики
памяти; останется только склад.
\end{intuition}

\section{Центральный тезис книги}

\begin{principle}[Главная формула всего курса]
Память Рин рассматривается как один типизированный темпоральный граф с
метаданными на узлах и рёбрах. Retrieval, обход, вывод, консолидация,
забывание, ночная переработка, настроение, личность и языковое выражение
трактуются как операторы над этим субстратом или как сцепленные с ним
процессы, а не как набор несогласованных подсистем с разными моделями мира.
\end{principle}

Именно этот тезис задаёт порядок изложения. Сначала вводится архитектурная
карта системы (глава 3). Затем даётся словарь основных объектов (глава 4).
После этого разбирается субстрат памяти (глава 5) и его многомасштабная
динамика (глава 11). Далее рассматриваются основные процессы: настроение
(глава 7), обход графа (глава 6), fusion и консолидация (глава 10), эпистемика
и ревизия убеждений (главы 13 и 14), влияние топологии (глава 8), ночная
работа (глава 12), слой личности и голоса (главы 16 и 17), инженерный рантайм
(глава 18) и, наконец, экспериментальная проверка (глава 22).

\section{Четыре процесса, которые нельзя настраивать по отдельности}

При первом знакомстве с задачей возникает соблазн мыслить модулями: сначала
``сделать retrieval'', потом ``добавить забывание'', потом ``приделать сон''.
Такой подход удобен управленчески, но почти всегда ошибочен содержательно.
В памяти Рин особенно тесно сцеплены четыре процесса.

\begin{enumerate}
\item Извлечение релевантного. Что именно всплывает в ответ на запрос, внутренний импульс или текущую сцену.
\item Пропагация уверенности. Как новый факт усиливает, ослабляет или ставит под сомнение соседние убеждения.
\item Забывание и консолидация. Что тускнеет, что переживает время, а что сворачивается в обобщение.
\item Ночная переработка. Как память в фоне инспектирует сама себя, чинит противоречия и перестраивает структуру.
\end{enumerate}

Эти процессы читают и пишут одни и те же поля одного и того же графа.
Изменение уверенности меняет то, что будет извлечено. Изменение топологии
меняет маршруты обхода. Консолидация меняет последующую пропагацию. Ночной
скраб меняет судьбу гипотез, созданных дневными операторами.

\begin{pitfall}[Ошибка независимой настройки]
Если каждый оператор проектируется как будто остальные существуют где-то
рядом, но не вмешиваются в его поведение, то система может хорошо работать в
локальном тесте и при этом распадаться в общем цикле. Например, ``удачный
retrieval'' на фиксированном графе может оказаться неудачным после первой же
волны консолидации; ``аккуратная пропагация'' может превратиться в лавину при
иной топологии; ``красивое забывание'' может уничтожить материал, без которого
ревизия убеждений станет невозможной.
\end{pitfall}

Именно поэтому завершающая часть книги посвящена не красивым примерам, а
экспериментальному стенду с контролируемой истиной (глава 22). Для такой
архитектуры особенно важно мерить не изолированные фрагменты поведения, а их
совместную физику.

\section{Чего эта книга не утверждает}

Текст следует читать как рабочую архитектурную теорию памяти для RIN, а не
как финальную общую теорию разума. В нём есть три разных слоя утверждений, и
их полезно различать с самого начала.

\begin{enumerate}
\item Твёрдые инженерные решения. Например, требование единого субстрата или различение активации и уверенности.
\item Рабочие гипотезы. Например, конкретные формы обхода, консолидации или эпистемических предохранителей.
\item Дальние линии развития. Например, переход к гиперграфам, клеточным комплексам или иным более выразительным субстратам.
\end{enumerate}

\begin{intuition}[Граница между моделью и реальностью]
Ни одна архитектура памяти не обязана быть истинной в сильном смысле слова;
она обязана быть содержательно плодотворной, инженерно управляемой и
проверяемой на стенде. Поэтому в этой книге важнее не вопрос ``похоже ли это
на окончательную истину о памяти'', а вопрос ``даёт ли такая схема хорошие
механизмы, понятные сбои и измеримые компромиссы''.
\end{intuition}

\section{Как читать этот курс}

Книга задумана так, чтобы не требовать от читателя постоянного угадывания,
``что именно автор имел в виду''. Ради этого соблюдаются два правила.

\begin{enumerate}
\item Новый объект по возможности сначала мотивируется, затем называется, затем формализуется.
\item Формула появляется там, где без неё уже становится неточно, а не там, где удобно сразу обозначить символы.
\end{enumerate}

Если нужен быстрый вход в систему, достаточно прочитать главы 3, 4, 5 и 20.
Если нужен полный контур, текст лучше читать по порядку: в этой версии книги
порядок глав соответствует логике постепенного усложнения модели.
